\section{Двигатель Евклида: невозможность вечного спуска}\label{sec:engine}

Программа сводит гипотезу о простых-близнецах к поведению одного динамического объекта --- цепи
самоподобных евклидовых уравнений, которую мы называем \emph{двигателем Евклида}, --- и опирает всё
здание на единственный факт: двигатель не может работать вечно. Предположим, что близнецов лишь
конечное число. Тогда, начиная с некоторого масштаба, каждая пара $(6m-1,\,6m+1)$ содержит составную
сторону, разложение которой обнажает активный делитель; его удаление порождает структурно такое же
евклидово уравнение строго меньшего центра. Повторяя, получаем бесконечную цепь чистых спусков, ни
на одном шаге не выходящую из clean-ядра, --- вечно работающий двигатель, совершающий нетривиальную
работу спуска и ничего при этом не теряющий. Мы доказываем, что такого двигателя не существует. Это
бесконечный спуск Ферма, переформулированный для центров евклидовых пар: состояние описывается одним
натуральным числом --- его \emph{высотой}, --- а один успешный шаг уменьшает высоту строго и
\emph{мультипликативно}, после чего работает лишь арифметика $\N$. Всё нижеследующее проверено машинно
на голом ядре Lean~4, без \code{mathlib} (\code{EuclidsPath/Engine/EPMI.lean},
\code{EuclidsPath/Engine/Irreversibility.lean}, \code{EuclidsPath/Engine/UniversalEngine.lean})~\cite{lean4,euclidspath-repo}.

\subsection{Шаг спуска и его строгость}

Абстрагируемся от арифметической начинки и оставим лишь то, что нужно для доказательства
невозможности. Состояние несёт единственную числовую характеристику --- высоту $h\in\N$,
содержательно центр $m$ пары $(6m-1,\,6m+1)$.

\begin{definition}[\code{DescentStep}]\label{thm:DescentStep}
Пусть $A\in\N$ --- фиксированный масштабный порог (нижняя граница активных множителей). Один успешный
clean-спуск с высоты $h$ на высоту $h'$ есть отношение
\begin{equation}\label{eq:descentstep}
  \mathrm{DescentStep}\,A\,h\,h' \;:=\; A\cdot h' < h.
\end{equation}
\end{definition}

Шаг фиксирует не абы какое, а \emph{мультипликативное} убывание: новая высота меньше в $A$ раз (с
точностью до строгости). Именно множитель $A$, а не разность, задаёт скорость спуска и делает
конечность неизбежной.

\begin{lemma}[\code{descent_strict}]\label{thm:descent_strict}
\green{}\; Если $1\le A$ и $\mathrm{DescentStep}\,A\,h\,h'$, то $h'<h$.
\end{lemma}

\emph{Идея доказательства.} Из $1\le A$ следует $h'\le A\cdot h'$ (умножение на положительный
множитель не уменьшает), а $A\cdot h'<h$ --- по \cref{eq:descentstep}; транзитивность даёт $h'<h$.
Лемма отделяет \emph{направление} спуска (высота падает) от его \emph{скорости} (падает в $A$ раз).

\subsection{Невозможность бесконечного спуска}

\begin{theorem}[\code{no_infinite_descent}]\label{thm:no_infinite_descent}
\green{}\; Пусть $1\le A$. Не существует последовательности высот $H:\N\to\N$, для которой каждый шаг
был бы успешным $A$-спуском, то есть
\begin{equation}\label{eq:perpetualheights}
  \forall\, t,\quad \mathrm{DescentStep}\,A\,(H\,t)\,(H\,(t+1)).
\end{equation}
Из существования такой последовательности выводится \code{False}.
\end{theorem}

\begin{proof}
Рассмотрим величину $H\,t+t$. На каждом шаге высота падает как минимум на $1$ (\cref{thm:descent_strict}),
а счётчик $t$ растёт на $1$; значит $H\,t+t$ не возрастает, и по индукции $\forall\,t,\ H\,t+t\le H\,0$.
Подставив $t=H\,0+1$, получаем $H\,(H\,0+1)+(H\,0+1)\le H\,0$, левая часть которого уже превосходит
$H\,0$ --- противоречие. В Lean индукция закрывается тактикой \code{omega}; \code{\#print axioms}
подтверждает, что теорема не зависит ни от каких аксиом и свободна от \code{sorry}.
\end{proof}

Мультипликативное убывание $H_t<H_0/A^{t}$ уводит высоту ниже $1$ за конечное число шагов, а
положительное целое не может быть меньше $1$: натуральный ряд имеет дно, и дискретный спуск обязан его
достичь. Для самого противоречия достаточно $A=1$ (порядковая полнота $\N$); множитель $A\ge 1$ даёт
количественную скорость. Это сильнее физического второго закона: поскольку высота --- целое со строгим
дном, спуск не просто затухает асимптотически, а останавливается не более чем за $H\,0$ шагов
(\cref{thm:turned_engine_halts}).

\subsection{Структурная форма: состояние и boundary-exit}

Чтобы связать абстрактную теорему с содержательной картиной, введём явное состояние и оператор шага,
различающий два исхода удаления активного множителя.

\begin{definition}[\code{State}]\label{thm:State}
Состояние двигателя есть структура с единственным полем $\mathrm{height}:\N$; содержательно это центр
$m$ текущей евклидовой пары. Шаг над состоянием $s$ --- индуктивный тип с двумя конструкторами,
\begin{equation}\label{eq:step}
  \mathrm{Step}\,A\,s \;=\;
  \begin{cases}
    \texttt{clean}\ s'\ (A\cdot s'.\mathrm{height} < s.\mathrm{height}), & \text{успешный спуск в clean-состояние,}\\[2pt]
    \texttt{boundary}, & \text{поглощающий выход } \bot,
  \end{cases}
\end{equation}
где ветвь \texttt{clean} несёт свидетельство мультипликативного убывания, а \texttt{boundary}
фиксирует выход из clean-ядра на границу, откуда нет возврата.
\end{definition}

\begin{theorem}[\code{no_perpetual_engine}]\label{thm:no_perpetual_engine}
\green{}\; Пусть $1\le A$. Нет траектории $\mathrm{run}:\N\to\mathrm{State}$, в которой каждый шаг ---
успешный clean-спуск, то есть
$A\cdot(\mathrm{run}\,(t+1)).\mathrm{height} < (\mathrm{run}\,t).\mathrm{height}$ для всех $t$.
\end{theorem}

\emph{Идея доказательства.} Приведение к абстрактной форме: берём высоты
$t\mapsto(\mathrm{run}\,t).\mathrm{height}$ и применяем \cref{thm:no_infinite_descent}. Это и есть
«нет вечного двигателя Евклида».

\begin{theorem}[\code{boundary_dichotomy}]\label{thm:boundary_dichotomy}
\green{}\; Для любого шага $st:\mathrm{Step}\,A\,s$ выполнено ровно одно из двух:
\begin{equation}\label{eq:dichotomy}
  \bigl(\exists\, s'\,h,\ st = \texttt{clean}\ s'\ h\bigr)\ \lor\ \bigl(st = \texttt{boundary}\bigr).
\end{equation}
\end{theorem}

\emph{Идея доказательства.} Разбор случаев по конструкторам \cref{eq:step}. Это типовая фиксация
boundary-exit law: descended-центр либо остаётся clean (высота строго упала в $A$ раз), либо выходит
на поглощающую границу $\bot$, откуда нет clean-продолжения той же ветви. Граница поглощает:
$\bot\not\to S$.

\subsection{Необратимость и стрела времени}

\cref{thm:no_infinite_descent} отвечает на вопрос «всегда ли двигатель останавливается?». Соседний
модуль \code{Irreversibility.lean} закрывает «может ли он повернуть назад?», давая дискретный аналог
второго закона термодинамики. Всё это доказано машинно только на стандартных аксиомах, без обращения
к распределению простых.

\begin{theorem}[\code{engine_never_returns}]\label{thm:engine_never_returns}
\green{}\; Пусть $1\le A$ и $H:\N\to\N$ удовлетворяет \cref{eq:perpetualheights}. Тогда $H$ строго
антимонотонна:
\begin{equation}\label{eq:strictanti}
  \mathrm{StrictAnti}\,H \;:\Longleftrightarrow\; \bigl(\forall\, s\,t,\; s<t \Rightarrow H\,t<H\,s\bigr).
\end{equation}
\end{theorem}

\emph{Идея доказательства.} Локальное строгое убывание $H(n+1)<H(n)$ из \cref{thm:descent_strict}
поднимается до глобальной антимонотонности через \code{strictAnti_nat_of_succ_lt}: на $\N$ убывание
на каждом единичном шаге эквивалентно убыванию на любом интервале. Высота есть координата
термодинамического времени, а \cref{eq:strictanti} --- формальная стрела времени: двигатель, тронувшись
в спуск, не может восстановить прежнее состояние. Это описывает любую цепь успешных спусков, \emph{буде
такая дана}; может ли она быть бесконечной --- отрицательно решает \cref{thm:no_infinite_descent}.

\begin{theorem}[\code{fuel_ascent_strictMono}]\label{thm:fuel_ascent_strictMono}
\green{}\; Отображение $n\mapsto n+1$ строго монотонно: $\mathrm{StrictMono}\,(\lambda n.\,n+1)$.
\end{theorem}

\emph{Идея доказательства.} Тривиально $n<n+1$. Тривиальность формы скрывает содержание: successor-цепь
$0,1,2,\dots$ --- бесконечная строго возрастающая траектория. Прибавление $+1$ --- «топливо»: бо́льших
центров всегда хватает. Сопоставление с \cref{thm:no_infinite_descent} обнажает асимметрию
термодинамической оси --- вверх бесконечно, вниз конечно. Эта асимметрия и есть стрела.

\begin{theorem}[\code{turned_engine_halts}]\label{thm:turned_engine_halts}
\green{}\; Пусть $H:\N\to\N$ и двигатель сделал $k$ строгих шагов вниз, то есть $H(t+1)<H(t)$ для всех
$t<k$. Тогда
\begin{equation}\label{eq:haltbound}
  k \;\le\; H(0).
\end{equation}
\end{theorem}

\emph{Идея доказательства.} Индукцией доказывается инвариант $H(t)+t\le H(0)$ для всех $t\le k$;
подстановка $t=k$ и отбрасывание $H(k)\ge 0$ дают \cref{eq:haltbound}. Это тот же баланс «высота плюс
время не растёт», что и в \cref{thm:no_infinite_descent}, прочитанный как конечная оценка, а не как
выход на противоречие: любой поворот вниз --- конечный путь длиной не более начальной высоты.

\subsection{Универсальная форма: двигатель против фундированности}

Невозможность выше есть тень единственного порядково-теоретического факта, формализованного в
\code{UniversalEngine.lean} над произвольным отношением. Это чтение несёт вся статья: где носитель
фундирован, двигатель запрещён \emph{безусловно}; где нет --- двигатель работает.

\begin{definition}[\code{PerpetualEngine}]\label{thm:PerpetualEngine}
Для отношения $r$ на типе $\alpha$ вечный двигатель есть бесконечная строго $r$-убывающая цепь:
\begin{equation}\label{eq:perpetualengine}
  \mathrm{PerpetualEngine}\,r \;:=\; \exists\, f:\N\to\alpha,\ \forall\, n,\ r\,(f(n+1))\,(f\,n).
\end{equation}
\end{definition}

\begin{theorem}[\code{perpetualEngine_iff_not_wellFounded}]\label{thm:perpetualEngine_iff_not_wellFounded}
\green{}\; Для любого $r$ на $\alpha$
\begin{equation}\label{eq:dividingline}
  \mathrm{PerpetualEngine}\,r \;\Longleftrightarrow\; \neg\,\mathrm{WellFounded}\,r.
\end{equation}
\end{theorem}

\emph{Идея доказательства.} Разворачивая $\mathrm{WellFounded}$ через характеризацию убывающими цепями,
двигатель \cref{eq:perpetualengine} есть буквально убывающая цепь, отсутствие которой и есть
фундированность. Это точная разделяющая линия: двигатель и фундированность взаимно дополнительны.

\begin{theorem}[\code{no_perpetual_engine_of_wellFounded}]\label{thm:no_perpetual_engine_of_wellFounded}
\green{}\; Если $\mathrm{WellFounded}\,r$, то $\neg\,\mathrm{PerpetualEngine}\,r$.
\end{theorem}

\emph{Идея доказательства.} Немедленно из \cref{thm:perpetualEngine_iff_not_wellFounded}. Невозможность
внутренняя, почти определительная: собственная определяющая цепь двигателя есть ровно то, что запрещает
фундированность. \cref{thm:no_infinite_descent} --- частный случай $r={<}$ на $\N$
(\code{no_perpetual_engine_on_nat}), с высотой в роли ранга.

\begin{theorem}[\code{no_perpetual_engine_of_rank}]\label{thm:no_perpetual_engine_of_rank}
\green{}\; Пусть $\rho:\alpha\to\beta$ с фундированным $(\beta,s)$, и каждый $r$-шаг строго уменьшает
ранг: $\forall\, x\,y,\ r\,x\,y \Rightarrow s\,(\rho\,x)\,(\rho\,y)$. Тогда
$\neg\,\mathrm{PerpetualEngine}\,r$.
\end{theorem}

\emph{Идея доказательства.} Двигатель $f$ для $r$ переносился бы вдоль $\rho$ в убывающую цепь
$\rho\circ f$ в фундированном $(\beta,s)$, противореча \cref{thm:no_perpetual_engine_of_wellFounded}.
$\N$-ранговая специализация (\code{no_perpetual_engine_of_natRank}) есть унификация: всякий фронт спуска
программы --- следствие этой одной теоремы, получаемое предъявлением ранга, строго убывающего вдоль
шагов.

\begin{theorem}[\code{perpetualEngine_on_real}]\label{thm:perpetualEngine_on_real}
\green{}\; На континууме двигатель работает: $\mathrm{PerpetualEngine}\,(<)$ на $\R$.
\end{theorem}

\emph{Идея доказательства.} Свидетель $(1/2)^n$ --- бесконечная строго $(<)$-убывающая цепь в $\R$;
значит, по \cref{thm:perpetualEngine_iff_not_wellFounded}, строгий порядок на $\R$ не фундирован. Это
честная граница контроллёра: на дискретных носителях двигатель запрещён и спуск обрывается, а на
вещественных вечный спуск реализуется. Запрет живёт в дискретности, а не в анализе.

\medskip

Мы ввели центральный объект и доказали, машинно и без аксиом, что бесконечно он работать не может
(\cref{thm:no_infinite_descent,thm:no_perpetual_engine}), не поворачивает назад
(\cref{thm:engine_never_returns}) и, свернув вниз, останавливается за конечное число шагов
(\cref{thm:turned_engine_halts}); дихотомия (\cref{thm:boundary_dichotomy}) оставляет поглощающую
границу единственной альтернативой продолжению. Честно очертим границу доказанного: это невозможность
бесконечного \emph{локального} clean-спуска вдоль одной ветви. Глобальный результат о непокрытии отсюда
\emph{не следует} автоматически --- мыслимо конечное дерево ограниченной глубины, каждая ветвь которого
честно обрывается absorbing boundary-leaf и которое при этом покрывает все центры. План закрытия
поднимает запрет с ветви на всё дерево, переоформленный как multi-rank non-cover с редукцией ранга к
rank-$1$ обструкции соседа; этот единственный открытый узел разбирается в \cref{sec:reduction}.
Двигатель же поставляет локальный закон конечности, на котором стоит вся редукция.
